\begin{center}[10.]\par\medskip\textit{Einfachste Ausdrücke für die Wirkungen galvanischer Ströme.}\end{center}
Die Fundamentalebene verläuft durch das wirkende Stromelement $AB$ und den Punkt $C$, auf den es wirkt.

Die komplexen Größen, die die Positionen $B,C$ relativ zu $A$ bezeichnen, seien $\beta,\gamma$. Ferner sei $r$ der Betrag von $\beta$. Ist schließlich auch bei $C$ bereits ein Strom mit dem Element $CD$ vorhanden, so sei $\gamma+\delta+i'\zeta$ die komplexe Größe, die die Position von $D$ relativ zu $A$ bezeichnet. Dann erhält man

I.\quad wenn in $C$ ein Strom fließt, für die Kraft, die $AB$ auf dessen materiellen Träger ausübt,
\[
\frac{\delta}{r}\quad\mathrm{Im.}\quad\frac{\beta}{\gamma}
\]

II.\quad wenn in $C$ kein Strom fließt, aber bei $C$ eine Bewegung stattfindet, die durch $\varepsilon$ in der durch $BC$ verlaufenden Ebene bezeichnet wird, für die elektromotorische Kraft,
\[
\frac{\beta}{r}\quad\mathrm R\quad\frac{\varepsilon}{\gamma}
\]
oder in übersichtlicherer Schreibweise
\SourceLeaf{621}\[
\begin{gathered}
G\quad\text{wirkendes Stromelement}\\[4pt]
\left.\begin{gathered}g\\m\end{gathered}\right\}\quad
\left.\begin{gathered}\text{vorhandenes Stromelement}\\\text{vorhandene Bewegung}\end{gathered}\right\}\quad
\begin{gathered}\text{am Punkt, auf den}\\\text{eingewirkt wird}\end{gathered}\\[4pt]
\left.\begin{gathered}\gamma\\\mu\end{gathered}\right\}\quad
\text{Kraft zur Erregung}\Bigg\}\quad
\begin{gathered}\text{eines Stroms}\\\text{einer Bewegung}\end{gathered}\\[4pt]
r\quad\text{Abstand als Betrag der komplexen Größe }l
\end{gathered}
\]
\[
\begin{aligned}
\text{1)}\qquad r\mu&=g\,.\,\mathrm{J.}\,\frac{G}{l}\\[4pt]
\text{2)}\qquad r\gamma&=G\,.\,\mathrm{R.}\,\frac{m}{l}
\end{aligned}
\]

\begin{center}[11.]\par\medskip\textit{Geradlinige Polygone.}\end{center}
Der Punkt, auf den eingewirkt wird, sei der Ursprung. Dann ist
\[
\text{I. Für}\qquad X=\int\frac{\mathrm dx}{r}
\]
der Beitrag der ersten Seite $PP'$
\[
\frac{x'-x}{r}\log\frac{\operatorname{cotg}\tfrac{1}{2}\theta'}{\operatorname{cotg}\tfrac{1}{2}\theta}
\]
wobei $\theta,\theta'$ die Winkel zwischen $PP'$ und $0P,0P'$ sind und $PP'$ in derselben Richtung genommen wird. Die Größe hinter dem Zeichen log.\ ist
\[
\begin{aligned}
&=\frac{r'+\dfrac{(x'-x)x'+(y'-y)y'+(z'-z)z'}{PP'}}{r+\dfrac{(x'-x)x+(y'-y)y+(z'-z)z}{PP'}}\\[6pt]
&=\frac{PP'\,.\,r'+r'r'-(r',PP')}{PP'\,.\,r-rr+(r,PP')}
\end{aligned}
\]

II.\quad Für $V$ ist der Beitrag des Winkels bei $P'$ der Unterschied zwischen dem Winkel der Ebenen $0PP',0P'P''$ und $180^0$. Setzt man den Flächeninhalt eines sphärischen Dreiecks mit den Seiten $a,b,c$ $=\omega$, so gilt
\[
\sin\tfrac{1}{2}\omega=\frac{6\text{ Pyramide }0PP'P''}{4rr'r'\cos\tfrac{1}{2}a\,.\,\cos\tfrac{1}{2}b\,.\,\cos\tfrac{1}{2}c}
\]
\SourceLeaf{622}
\begin{center}[12.]\end{center}
\noindent$g$\quad Stärke eines galvanischen Stroms\par
\noindent$\mu$\quad Dichte des Magnetismus auf einer vom Rheophor begrenzten Fläche\par
\noindent$\rho$\quad Abstand dieser Fläche von einer zweiten, negativ magnetisierten
\[
\text{1.}\qquad g=\mu\rho
\]

2.\quad Wirkung eines elektrischen Elements auf ein anderes. Die Position des ersten relativ zum zweiten wird durch die komplexe Größe $u$ zur Zeit $t$ bestimmt, bei einem Abstand $=r$.
\[
-\frac{u}{r^3}-\frac{\alpha}{ur}\left(\frac{\mathrm du}{\mathrm dt}\right)^2-\frac{\beta}{r}\frac{\mathrm d\mathrm du}{\mathrm dt^2}
\]

3.\quad Gegenseitige Wirkung zwischen einem elektrischen und einem magnetischen Element bei der relativen Geschwindigkeit $v$
\[
\frac{\gamma v}{rr}
\]

4.\quad In einem galvanischen Strom der Stärke $g$ bewegt sich in der Zeit $t$ durch jeden Querschnitt die positive Elektrizität $\varepsilon gt$ in die eine Richtung und die negative $-\varepsilon gt$ in die andere.

5.\quad Es geht darum, die Beziehungen zwischen $\alpha,\beta,\gamma,\varepsilon$ zu bestimmen. Die Beziehung $\gamma=2\beta\varepsilon$ folgt aus der Induktion bei der Entstehung eines Stroms während der Zeit $\theta$, wobei die Kraft, solange der Strom entsteht —
\[
\begin{gathered}
=\frac{\gamma\pi hhg}{\theta RR}=\frac{2\beta\pi hh\varepsilon g}{\theta RR}\quad\text{wird}\\[6pt]
\gamma=4\alpha\varepsilon,\qquad\frac{\gamma\pi hhgv}{R^3}=\frac{4\alpha\pi hh\varepsilon gv}{R^3}\\[6pt]
2\gamma\varepsilon=1,\qquad\alpha=\frac{1}{8\varepsilon\varepsilon}=\tfrac{1}{2}\gamma\gamma,\qquad\beta=\frac{1}{4\varepsilon\varepsilon}=\gamma\gamma
\end{gathered}
\]

